\documentclass[../../main.tex]{subfiles}

\begin{document}

\section{弱序列紧性}\label{section:ch19-weak-compactness}

令 $(X,\|\cdot\|)$ 为赋范线性空间. 序列 $\{f_{n}\}$ 称为弱收敛于 $f\in X$, 若对每个有界线性泛函 $S\in X^{*}$, $\lim_{n\to\infty}S(f_{n})=S(f)$. 若 $1<p<\infty$ 且 $q$ 是共轭指数, 则根据 Riesz 表示定理, $L^{p}(X,\mu)$ 中的序列 $\{f_{n}\}$ 弱收敛于 $f$ 当且仅当对每个 $g\in L^{q}(X,\mu)$,
\begin{align*}
\lim_{n\to\infty}\int_{X}f_{n}g\,d\mu=\int_{X}fg\,d\mu.
\end{align*}

\begin{theorem}[自反性]\label{theorem:ch19-reflexivity}
令 $(X,\mathcal{M},\mu)$ 为 $\sigma$ 有限测度空间, 而 $1<p<\infty$. 则 Banach 空间 $L^{p}(X,\mu)$ 是自反的, 即典范嵌入 $J:L^{p}\to(L^{p})^{**}$ 是满射.
\end{theorem}

\begin{proof}
根据 Riesz 表示定理, $(L^{p})^{*}=L^{q}$ 且 $(L^{q})^{*}=L^{p}$, 其中 $1/p+1/q=1$. 因此 $J:L^{p}\to(L^{p})^{**}=(L^{q})^{*}=L^{p}$ 是恒等映射, 从而是满射.
\end{proof}

根据 Banach 空间的一般理论, 自反 Banach 空间的闭单位球是弱序列紧的: 每个有界序列有弱收敛的子序列.

\begin{theorem}[弱序列紧性]\label{theorem:ch19-weak-sequential-compactness}
令 $(X,\mathcal{M},\mu)$ 为 $\sigma$ 有限测度空间, $1<p<\infty$, 而 $\{f_{n}\}$ 是 $L^{p}(X,\mu)$ 中的有界序列: $\sup_{n}\|f_{n}\|_{p}<\infty$. 则存在子序列 $\{f_{n_{k}}\}$ 和 $f\in L^{p}(X,\mu)$, 使得 $\{f_{n_{k}}\}$ 弱收敛于 $f$.
\end{theorem}

\begin{proof}
由于 $L^{p}(X,\mu)$ 是自反的, 其闭单位球是弱紧的. 根据 Eberlein-\v{S}mulian 定理, 弱紧性等价于弱序列紧性.
\end{proof}

\begin{exercise}\label{exercise:ch19-10}
令 $1<p<\infty$. 证明: 若 $\{f_{n}\}$ 在 $L^{p}$ 中弱收敛于 $f$ 且 $\|f_{n}\|_{p}\to\|f\|_{p}$, 则在 $L^{p}$ 中 $f_{n}\to f$(范数收敛).
\end{exercise}

\end{document}